\documentclass{report}
\usepackage{amsmath,amssymb}
\setlength{\parindent}{0mm}
\setlength{\parskip}{1em}
\begin{document}
\begin{center}
\rule{6in}{1pt} \
{\large Matthew T. Calef \\
{\bf Performance and examination of a Krylov accelerator}}

PO Box 1663 \\ Mail Stop D413 \\ Los Alamos \\ NM 87545
\\
{\tt mcalef@lanl.gov}\\
Erin D. Fichtl\\
James S. Warsa\\
Markus Berndt\\
Neil N. Carlson\end{center}

We present a variant of Anderson Mixing, and compare it to the Broyden
and JFNK methods used to solve a discretization of the Boltzmann
equation. We show that, for the problems under consideration, our variant
of Anderson Mixing finds the solution in the fewest number of iterations.
In addition we examine and extend some recent theoretical work regarding
Anderson Mixing.


\end{document}
